% TEMPLATE-MILC-v3.tex - plantilla maestra UNICA de las guias MILC v3.
%
% La usan las cuatro familias de guias, cada una con su driver:
%   · Grados 6-11          -> scripts/build-guias-g11.py
%   · Bebras               -> scripts/build-guias-bebras.py
%   · Semillero            -> scripts/build-guias-semillero.py
%   · Territorio interior  -> scripts/build-guias-territorio-interior.py
%
% Antes habia tres copias casi identicas (~400 lineas iguales, 6 distintas):
% cada mejora habia que aplicarla tres veces, y en la practica no se hacia
% (el motor de recursos vivia solo en dos de ellas). Lo que varia entre
% familias esta parametrizado en 6 placeholders que cada driver llena
% (se nombran aqui SIN su sintaxis real: el builder reemplaza por texto plano,
% tambien dentro de los comentarios, y un valor multilinea rompe el preambulo):
%
%   HEADER_IZQ        encabezado izquierdo de pagina
%   HEADER_DER        encabezado derecho de pagina
%   PORTADA_ETIQUETA  rotulo sobre el numero grande (GRADO/MOMENTO/MODULO)
%   PORTADA_SERIE     linea de serie de la portada
%   PORTADA_ID        identificador de la guia en la portada
%   PORTADA_CIRCULO   el circulito de grado (°), o vacio si no aplica
%   PDF_TITULO        titulo en texto plano (sin \textbf ni comillas TeX) para
%                     los metadatos del PDF (/Title); lo produce lib_xelatex.texto_pdf
%
% Composicion (auditoria 2026-09-03): las bandas de estacion piden espacio
% (\Needspace*) y prohiben el salto justo despues (\nopagebreak), asi no quedan
% huerfanas al pie; las cajas largas (softbox, drawbox, infoband, puentebox) son
% `breakable`, asi una caja de media pagina ya no arrastra todo a la siguiente
% dejando la anterior al 40 %. Todo texto sobre mostaza o verde va en milcNegro
% (blanco sobre #E5B400 da 1,93:1 y sobre #58A923 2,95:1; ninguno pasa WCAG AA).
% El resto de claves las documenta content/guias/_SCHEMA.md. El script valida
% que no queden placeholders sin reemplazar antes de escribir el archivo final.

% Etiquetado PDF (latex-lab): /Lang, estructura y texto alternativo de las
% figuras, para lectores de pantalla. Debe ir ANTES de \documentclass.
\DocumentMetadata{lang=es-CO,tagging=on}
\documentclass[11pt,letterpaper]{article}
% headheight=24pt: la cabecera derecha lleva dos lineas (identidad de la guia
% y estacion actual). headsep se acorta en la misma medida para que el bloque
% de texto quede exactamente donde estaba con headheight=12pt.
\usepackage[margin=2.0cm,headheight=24pt,headsep=13pt]{geometry}
\usepackage[table]{xcolor}
\usepackage{fontspec}
\usepackage[spanish]{babel}
\usepackage{tikz}
% Librerías TikZ para los diagramas de las guías (bloque `recursos:`):
%  · babel  → babel en español vuelve activos caracteres como «>», y TikZ los usa
%             en sus opciones (>=stealth, ->). Sin esta librería, cualquier
%             diagrama con flechas revienta con «\language@active@arg> has an
%             extra }». Debe ir ANTES de que se dibuje nada.
%  · positioning        → sintaxis «below left=2pt and 3pt».
%  · decorations.pathreplacing → llaves ({brace}) para señalar rangos.
\usetikzlibrary{babel,positioning,decorations.pathreplacing}
\usepackage{graphicx}
% Circuitos: el trabajo de electrónica se hace hoy con micro:bit y sus sensores
% integrados (MakeCode, sin cableado), así que no se necesitan esquemáticos.
% Si algún día entran montajes con protoboard, basta descomentar esta línea:
% los diagramas .tex/.tikz declarados en `recursos:` ya se incrustan solos.
% \usepackage{circuitikz}
\usepackage[most]{tcolorbox}
\usepackage{tabularx}
\usepackage{array}
\usepackage{fancyhdr}
\usepackage{titlesec}
\usepackage[document]{ragged2e}  % bandera a la derecha en todo el cuerpo
\usepackage{enumitem}
\usepackage{microtype}
\usepackage{etoolbox}
\usepackage{needspace}
% hyperref al final del preambulo: metadatos del PDF (titulo, autor, idioma,
% familia), marcadores por estacion (\pdfbookmark) y enlaces sin recuadro.
\usepackage[hidelinks,
  pdftitle={Limpieza de datos — el oficio del archivero},
  pdfauthor={PhD. Álvaro Cárdenas Orozco},
  pdfsubject={Tecnología e Informática · Grado 9 · Guía 27 · MILC v3},
  pdflang={es-CO},
  bookmarksopen=true,bookmarksnumbered=false]{hyperref}

% Helvetica Neue no trae la flecha U+2192, asi que cada «→» escrito literalmente
% en el YAML se compilaba a NADA, en silencio (462 flechas en 61 guias: rutas del
% tipo «paleta → tipografia → lockup» salian rotas). Mapeamos el caracter a su
% version matematica, que si tiene glifo. Asi el autor puede seguir escribiendo
% «→» comodamente en el YAML.
\usepackage{newunicodechar}
\newunicodechar{→}{\ensuremath{\rightarrow}}
% Mismo problema con los simbolos de marca y vinetas: el «✓» de «una palomita ✓»
% salia como cuadrito vacio en el PDF de todas las guias que lo usaban.
\usepackage{pifont}
\newunicodechar{✓}{\ding{51}}
\newunicodechar{✗}{\ding{55}}
\newunicodechar{✔}{\ding{52}}
\newunicodechar{•}{\textbullet}
\newunicodechar{≥}{\ensuremath{\geq}}
\newunicodechar{≤}{\ensuremath{\leq}}

\setmainfont{Helvetica Neue}
\setsansfont{Helvetica Neue}
\setlength{\parindent}{0pt}
\setlength{\parskip}{6pt}
% Sin particion de palabras: en 6.º-7.º una palabra cortada por guion al final
% de linea («Comuni-cacion») frena la lectura mucho mas que una linea corta.
\hyphenpenalty=10000
\exhyphenpenalty=10000
\pretolerance=2000
\tolerance=3000
\setlength{\emergencystretch}{3em}
\linespread{1.10}
\renewcommand{\arraystretch}{1.25}
\setlist{leftmargin=5mm,itemsep=1mm,topsep=1mm}
\newcolumntype{Y}{>{\RaggedRight\arraybackslash}X}

\definecolor{milcMagenta}{HTML}{E6007E}
\definecolor{milcVino}{HTML}{5A0038}
\definecolor{milcTurquesa}{HTML}{0093A5}
\definecolor{milcVerde}{HTML}{58A923}
\definecolor{milcMostaza}{HTML}{E5B400}
\definecolor{milcOcre}{HTML}{B86600}
\definecolor{milcNegro}{HTML}{241816}
\definecolor{milcGris}{HTML}{F7F7F4}
\definecolor{milcLinea}{HTML}{E8E2DD}
\definecolor{milcGradePrimary}{HTML}{7C3AED}
\definecolor{milcGradeDark}{HTML}{4C1D95}
\definecolor{milcGradeSoft}{HTML}{EDE9FE}
\definecolor{milcGradeAccent}{HTML}{CFE7FF}
\definecolor{milcGradeText}{HTML}{FFFFFF}
\color{milcNegro}

\pagestyle{fancy}
\fancyhf{}
% Cabecera de dos lineas: identidad de la guia arriba y, a la derecha y en
% gris, la estacion en curso (\markright desde \milcestacion). Antes de la
% primera estacion la segunda linea va vacia; el \strut mantiene la altura
% para que la linea de arriba no baile entre paginas.
\lhead{\small Tecnología e Informática · Grado 9\\[-1pt]{\footnotesize\strut}}
\rhead{\small Guía 27 · MILC v3\\[-1pt]{\footnotesize\strut\color{milcNegro!65}\rightmark}}
\cfoot{\small\thepage}
\renewcommand{\headrulewidth}{0pt}

% Sobre mostaza (#E5B400) y verde (#58A923) el texto va en milcNegro; sobre el
% resto de la paleta, en blanco. Se usa dentro de fonttitle/fontupper, donde un
% \color posterior manda sobre coltitle.
\newcommand{\milctintasobre}[1]{%
  \ifboolexpr{test{\ifstrequal{#1}{milcMostaza}} or test{\ifstrequal{#1}{milcVerde}}}%
    {\color{milcNegro}}{\color{white}}}

\newtcolorbox{titlebox}[1]{
  enhanced,colback=#1,colframe=#1,arc=7mm,boxrule=0pt,
  left=7mm,right=7mm,top=3mm,bottom=3mm,
  before skip=8pt,after skip=8pt,
  fontupper=\sffamily\bfseries\Large\color{white}
}

\newtcolorbox{softbox}[3]{
  enhanced,breakable,pad at break=3mm,
  colback=#2,colframe=#1,borderline west={4pt}{0pt}{#1},
  arc=2mm,boxrule=.4pt,left=4mm,right=4mm,top=3mm,bottom=3mm,
  before skip=6pt,after skip=8pt,title={#3},coltitle=white,
  colbacktitle=#1,fonttitle=\sffamily\bfseries\milctintasobre{#1},fontupper=\RaggedRight,
  attach boxed title to top left={xshift=3mm,yshift=-2mm},
  boxed title style={arc=2mm, boxrule=0pt}
}

\newtcolorbox{drawbox}[2]{
  enhanced,breakable,pad at break=4mm,
  colback=white,colframe=#1,arc=4mm,boxrule=1pt,
  left=5mm,right=5mm,top=4mm,bottom=4mm,before skip=8pt,after skip=8pt,
  title={#2},coltitle=white,colbacktitle=#1,fonttitle=\sffamily\bfseries\milctintasobre{#1},
  fontupper=\RaggedRight,
  attach boxed title to top left={xshift=4mm,yshift=-2mm},
  boxed title style={arc=3mm, boxrule=0pt}
}

\newtcolorbox{infoband}[1]{
  enhanced,breakable,pad at break=2mm,colback=#1!8,colframe=#1!50,arc=2mm,boxrule=.5pt,
  left=4mm,right=4mm,top=2mm,bottom=2mm,before skip=4pt,after skip=4pt,
  fontupper=\small\RaggedRight
}

% Puente narrativo entre secciones (contrato editorial Conéctate).
% Da continuidad: "Acabas de X. Ahora vas a Y porque Z."
% Visualmente: banda izquierda en color del grado, texto en itálica pequeña.
\newtcolorbox{puentebox}{
  enhanced,breakable,pad at break=2mm,colback=milcGradeSoft,colframe=milcGradeDark,
  borderline west={3pt}{0pt}{milcGradeDark},
  arc=0mm,boxrule=0pt,left=5mm,right=4mm,top=2.5mm,bottom=2.5mm,
  before skip=4pt,after skip=8pt,
  fontupper=\itshape\small\color{milcNegro}\RaggedRight
}
% La flecha va en modo matematico a proposito: Helvetica Neue NO tiene el glifo
% U+2192, asi que \textrightarrow se compilaba a NADA («Missing character: There
% is no → in font Helvetica Neue») y el puente salia sin flecha. En modo
% matematico la flecha viene de la fuente matematica y si se dibuja.
\newcommand{\puente}[1]{%
  \begin{puentebox}%
    \textcolor{milcGradeDark}{$\rightarrow$}\hspace{2mm}#1%
  \end{puentebox}%
}

\newcommand{\linead}[1]{\noindent\color{milcLinea}\rule{#1}{0.5pt}\color{milcNegro}}
\newcommand{\respuesta}[1]{\par\vspace{2mm}\foreach \n in {1,...,#1}{\noindent\color{milcLinea}\rule{\linewidth}{0.5pt}\par\vspace{5mm}}\color{milcNegro}}
\newcommand{\checkbox}{\textcolor{milcGradeDark}{\fbox{\phantom{x}}}\hspace{5pt}}

% ─── Listas aireadas para las guías socioemocionales ────────────────────
% Componer las verificaciones como «\checkbox texto\\» deja el texto que salta
% de línea debajo del cuadrito y apelmaza el bloque. Estos entornos alinean la
% continuación y airean los ítems. Los usa Territorio Interior; cualquier
% familia puede hacerlo.
\newlist{checklista}{itemize}{1}
\setlist[checklista]{
  label=\textcolor{milcGradeDark}{\fbox{\phantom{x}}},
  leftmargin=9mm, labelsep=3.2mm, itemsep=2.4mm,
  topsep=2.5mm, parsep=0pt, partopsep=0pt, before=\RaggedRight
}
\newlist{pasoslista}{enumerate}{1}
\setlist[pasoslista]{
  label=\textcolor{milcGradeDark}{\sffamily\bfseries\arabic*.},
  leftmargin=8mm, labelsep=3mm, itemsep=2.8mm,
  topsep=1mm, parsep=0pt, partopsep=0pt, before=\RaggedRight
}

\newcommand{\phasecard}[4]{
\begin{tcolorbox}[enhanced,colback=#3,colframe=#2,arc=3mm,boxrule=.5pt,height=3.05cm,valign=center,halign=center,left=1.5mm,right=1.5mm,top=2mm,bottom=2mm]
{\sffamily\bfseries\fontsize{22}{24}\selectfont\textcolor{#2}{#1}}\par
{\sffamily\bfseries\small #4}
\end{tcolorbox}
}

% ── Piezas del rediseno 2026 (legibilidad 6.º-11.º) ───────────────────
% Banda de estacion: reemplaza el titlebox macizo. Titulo a la izquierda,
% dato practico (tiempo/modalidad) a la derecha, en una sola linea.
% Nunca huerfana: pide 9 lineas libres (o salta de pagina rellenando con
% \vfill) y prohibe el salto justo despues de la banda. Nueve y no seis:
% la caja partible que sigue necesita ~80pt (titulo adosado, relleno y dos
% lineas) para arrancar en la misma pagina; con menos, tcolorbox la manda
% entera a la siguiente y la banda se queda sola. El penalty 10000 va
% en `after`, ANTES del espacio posterior: un `after skip` (aunque sea 0pt)
% es un \vskip tras la caja, y ese glue es un punto de corte legal que un
% \nopagebreak escrito despues de \end{tcolorbox} ya no cierra. Cada banda
% deja marcador PDF y actualiza la cabecera derecha.
\newcounter{milcestacion}
\newcommand{\milcestacion}[2]{%
\Needspace*{9\baselineskip}%
\stepcounter{milcestacion}%
\pdfbookmark[1]{#2}{milcestacion\themilcestacion}%
\markright{#2}%
\begin{tcolorbox}[enhanced,colback=#1,colframe=#1,arc=2mm,boxrule=0pt,
  left=5mm,right=5mm,top=2.6mm,bottom=2.6mm,before skip=11pt,
  after={\par\nopagebreak\vskip7pt}]
{\sffamily\bfseries\fontsize{15}{18}\selectfont\milctintasobre{#1}#2}
\end{tcolorbox}}

% Tarjeta de paso numerada: sustituye las tablas «Pilar / Aplicacion», que
% eran formato de consulta docente, no de estudiante. El numero va en un
% circulo sobre el borde para que la secuencia se lea de un vistazo.
\newcommand{\milcpaso}[3]{%
\Needspace{4\baselineskip}%
\begin{tcolorbox}[enhanced,colback=white,colframe=milcLinea,arc=1.5mm,boxrule=.9pt,
  left=14mm,right=4mm,top=3mm,bottom=3mm,before skip=4pt,after skip=4pt,
  fontupper=\RaggedRight,
  overlay={\node[anchor=north west,circle,fill=#2,text=white,inner sep=0pt,
    minimum size=7.4mm,font=\sffamily\bfseries\small]
    at ([xshift=3.6mm,yshift=-3.2mm]frame.north west) {#1};}]
#3
\end{tcolorbox}}

% Casilla marcable de tamano util con lapiz (la anterior era \fbox{\phantom{x}},
% de alto variable segun el texto que la seguia).
\newcommand{\milcmarca}{\framebox[3.8mm]{\rule{0pt}{3.4mm}}}

% Fila marcable de lista suelta (checks de la estacion 1).
\newcommand{\milccheckrow}[1]{%
\par\vspace{1.8mm}\noindent
\makebox[6.4mm][l]{\raisebox{-.55mm}{\textcolor{milcNegro}{\milcmarca}}}%
\parbox[t]{\dimexpr\linewidth-6.4mm\relax}{\RaggedRight #1}\par}

% Celda marcable dentro de la rubrica (tabla de autoevaluacion). Misma
% sangria francesa, pero pensada para vivir dentro de una columna X.
\newcommand{\milcopcion}[1]{%
\makebox[5.2mm][l]{\raisebox{-.55mm}{\textcolor{milcNegro}{\milcmarca}}}%
\parbox[t]{\dimexpr\linewidth-5.2mm\relax}{\RaggedRight #1}}

% ── Recursos visuales de la guía ──────────────────────────────────────
% Declarados en el bloque `recursos:` del YAML y emitidos por el builder.
% \guiaFigura   → imagen rasterizada o PDF (foto, carta celeste, montaje).
% \guiaDiagrama → fuente TikZ/circuitikz (.tex) incrustada como vector:
%                 es la vía para circuitos y esquemáticos editables.
% [#1] = texto alternativo (alt) · #2 = ruta relativa al .tex · #3 = pie
% El alt llega al PDF etiquetado (/Alt) y lo lee el lector de pantalla.
\newcommand{\guiaFigura}[3][]{%
\begin{center}
\includegraphics[width=0.88\linewidth,height=0.38\textheight,keepaspectratio,alt={#1}]{#2}\\[1.5mm]
{\footnotesize\itshape\textcolor{milcNegro!75}{#3}}
\end{center}
}

\newcommand{\guiaDiagrama}[2]{%
\begin{center}
\input{#1}\\[1.5mm]
{\footnotesize\itshape\textcolor{milcNegro!75}{#2}}
\end{center}
}

\begin{document}
\thispagestyle{empty}

% ── Portada ───────────────────────────────────────────────────────────
% Antes: un rectangulo de color a pagina completa. Se veia bien en pantalla,
% pero en la fotocopiadora del colegio se comia el toner y salia gris sucio.
% Ahora la banda ocupa el tercio superior y el resto va en blanco.
\begin{tikzpicture}[remember picture,overlay]
  \path[fill=milcGradePrimary]
    (current page.north west) rectangle ([yshift=-10.6cm]current page.north east);

  \node[anchor=north west,text=milcGradeText] at ([xshift=2.0cm,yshift=-1.55cm]current page.north west)
    {\sffamily\bfseries\fontsize{12}{14}\selectfont GRADO};
  \node[anchor=north west,text=milcGradeText,opacity=.85] at ([xshift=2.0cm,yshift=-2.35cm]current page.north west)
    {\sffamily\bfseries\fontsize{8.5}{10}\selectfont SERIE GUÍAS · TECNOLOGÍA E INFORMÁTICA};
  \node[anchor=north east,text=milcGradeText,opacity=.85,align=right] at ([xshift=-2.0cm,yshift=-1.55cm]current page.north east)
    {\sffamily\bfseries\fontsize{9}{12}\selectfont I.E. Sor María Juliana\\Cartago, Valle del Cauca};

  \node[anchor=west,text=milcGradeText] (gradonum) at ([xshift=1.85cm,yshift=-7.1cm]current page.north west)
    {\sffamily\bfseries\fontsize{112}{112}\selectfont 9};
  % El circulito de grado (°) solo aplica a las guias de grado («11°»); en
  % Bebras y Semillero el numero grande es un momento/modulo, no un grado.
  % Se posiciona relativo al nodo `gradonum`, no en coordenadas absolutas.
  \draw[milcGradeText,line width=1.6pt]
    ([xshift=2.0mm,yshift=-4.5mm]gradonum.north east) circle (.15cm);

  \node[anchor=west,text=milcGradeText,text width=11.4cm,align=left] at ([xshift=1.5cm,yshift=0.5cm]gradonum.east)
    {
     {\sffamily\bfseries\fontsize{9}{11}\selectfont\MakeUppercase{Período 3 · Datos --- del registro al insight}}\\[3mm]
     {\sffamily\bfseries\fontsize{21}{25}\selectfont\setlength{\baselineskip}{27pt}Limpieza de datos ---\\[4pt]el oficio\\[4pt]del archivero}\\[3.5mm]
     {\sffamily\bfseries\fontsize{15}{18}\selectfont Guía 27-9-TIC}
    };
\end{tikzpicture}

\vspace*{9.4cm}
\vfill

{\sffamily\bfseries\fontsize{8}{10}\selectfont\textcolor{milcGradeDark}{LO QUE TE LLEVAS HOY}}

\begin{tcolorbox}[enhanced,colback=white,colframe=milcNegro,arc=2mm,boxrule=1.6pt,
  left=5mm,right=5mm,top=4mm,bottom=4mm,before skip=3pt,after skip=12pt,
  fontupper=\RaggedRight\fontsize{11}{15.5}\selectfont]
Un conjunto de datos sucio, limpiado sobre una copia y con el original intacto: duplicados resueltos, formatos normalizados, vacíos decididos y errores de escritura corregidos. Lo acompaña la bitácora con al menos cuatro entradas, cada una diciendo qué encontraste, qué decidiste, por qué y a qué celdas afectó. La prueba: otra persona podría repetir tu limpieza leyendo solo la bitácora.
\end{tcolorbox}

\begin{tcolorbox}[enhanced,colback=milcGris,colframe=milcGris,arc=2mm,boxrule=0pt,
  left=5mm,right=5mm,top=3.5mm,bottom=3.5mm,after skip=12pt]
{\sffamily\bfseries\fontsize{8}{10}\selectfont\textcolor{milcNegro!55}{ESTA GUÍA ES DE}}\\[3mm]
\begin{tabularx}{\linewidth}{X p{3.2cm} p{3.2cm}}
\linead{\linewidth} & \linead{3.0cm} & \linead{3.0cm}\\[-1mm]
{\fontsize{8.5}{10}\selectfont\textcolor{milcNegro!55}{Nombre completo}} &
{\fontsize{8.5}{10}\selectfont\textcolor{milcNegro!55}{Grupo}} &
{\fontsize{8.5}{10}\selectfont\textcolor{milcNegro!55}{Fecha}}\\
\end{tabularx}
\end{tcolorbox}

\vfill

\noindent\textcolor{milcLinea}{\rule{\linewidth}{1pt}}\\[2mm]
{\sffamily\bfseries\fontsize{9.5}{12}\selectfont Autor: Dr. Álvaro Cárdenas Orozco}\\[1mm]
{\sffamily\fontsize{8.5}{11}\selectfont\textcolor{milcNegro!55}{Metodología MILC · apertura ancestral · pensamiento computacional · triángulo Dussel–estoicismo–Floridi}}

\newpage

\section*{Apertura: Limpieza de datos --- el oficio del archivero}

\begin{softbox}{milcGradeDark}{milcGradeSoft}{Saber ancestral}
Cuando una comunidad indígena va a constituir o ampliar un resguardo, hay que levantar un \textbf{censo poblacional}. Conviene decir qué es exactamente: un procedimiento del Estado hecho con la comunidad, no una práctica ancestral. El instructivo lo define como una fotografía de la realidad de una comunidad en un tiempo y un espacio específicos. Se levanta persona por persona, con apoyo de miembros del cabildo, profesores y líderes, y después funcionarios lo verifican. Y aquí está lo que importa para hoy: la tierra se distribuye conforme al censo realizado. Eso convierte una lista en algo muy serio. La \textbf{cara de exclusión} cabe en una frase: quien no queda en la lista el día del censo no existe para el reparto. Un dato mal escrito, un nombre repetido o una casilla vacía no son errores de forma. Son alguien que aparece dos veces o alguien que no aparece.\par\smallskip{\footnotesize\textcolor{milcNegro!70}{\textbf{Fuente.} Agencia Nacional de Tierras, Dirección de Asuntos Étnicos. (2023). \emph{Instructivo: levantamiento y diligenciamiento del formato de censo poblacional para comunidades indígenas} (ACCTI-I-021, versión 1).}}
\end{softbox}

\begin{softbox}{milcTurquesa}{milcGris}{Saber contemporáneo}
Un conjunto de datos se ensucia de cuatro maneras y cada una tiene su tratamiento. \textbf{Duplicados}: filas que repiten información. Se elimina la copia, se conserva el original y se anota. \textbf{Formatos inconsistentes}: la misma fecha escrita de tres modos, o mayúsculas mezcladas. Se elige un formato y se normaliza todo. \textbf{Vacíos}: se decide si se completan, se marcan con un valor explícito o se descarta la fila, y se dice cuál de las tres. \textbf{Errores de escritura}: nombres mal tecleados que parten en dos lo que era uno. El método es siempre el mismo: identificar, decidir y documentar, en ese orden. Nunca se corrige sin haber decidido, y nunca se decide sin dejar constancia. Y hay una regla que no se negocia: el original se conserva, la copia se limpia. Lo que se borra sin registro no se puede recuperar ni discutir.
\end{softbox}

\begin{softbox}{milcMostaza}{milcGris}{Pregunta-puente}
\textit{En un censo, quien queda fuera de la lista no existe para el reparto. En tu conjunto de datos, ¿qué desaparece cada vez que borras una fila sin anotarlo?}
\end{softbox}

\begin{softbox}{milcVerde}{milcGris}{Saber y saber hacer}
Al terminar podrás: (1) \textbf{identificar} los cuatro tipos de suciedad en un conjunto de datos que no habías visto antes; (2) \textbf{explicar} qué tratamiento pide cada uno y por qué la decisión va antes que la corrección; (3) \textbf{aplicar} la limpieza sobre una copia y documentarla para que otro pueda repetirla.
\end{softbox}



\small
\begin{tabularx}{\textwidth}{>{\bfseries}p{3.0cm}Y}
\rowcolor{milcGradePrimary!12}Autor & Dr. Álvaro Cárdenas Orozco\\
DBA & Tratamiento y visualización honesta de datos para la toma de decisiones (MEN, T\&I)\\
Referentes & Ética del dato · Visualización honesta · Pensamiento computacional\\
Producto final & Un conjunto de datos sucio, limpiado sobre una copia y con el original intacto: duplicados resueltos, formatos normalizados, vacíos decididos y errores de escritura corregidos. Lo acompaña la bitácora con al menos cuatro entradas, cada una diciendo qué encontraste, qué decidiste, por qué y a qué celdas afectó. La prueba: otra persona podría repetir tu limpieza leyendo solo la bitácora.\\
\end{tabularx}
\normalsize

\puente{En las sesiones anteriores recogiste datos, los pusiste en forma, calculaste y graficaste. Hoy trabajas la parte que se lleva la mayor parte del tiempo de cualquiera que trabaje con datos y que casi nunca se enseña. En la sesión 8 miras gráficos que engañan.}

\milcestacion{milcGradeDark}{Ruta MILC: así vamos a aprender}
\begin{tabularx}{\textwidth}{>{\centering\arraybackslash}X>{\centering\arraybackslash}X>{\centering\arraybackslash}X>{\centering\arraybackslash}X}
\phasecard{1}{milcVerde}{milcGris}{Escucha\\[-1mm]\small Observo, escucho y pregunto.} &
\phasecard{2}{milcTurquesa}{milcGris}{Entender\\[-1mm]\small Sistematizo y ordeno.} &
\phasecard{3}{milcMagenta}{milcGris}{Hacer\\[-1mm]\small Construyo y aplico.} &
\phasecard{4}{milcVino}{milcGris}{Cierre\\[-1mm]\small Evalúo y reflexiono.}
\end{tabularx}


\puente{Antes de la teoría vas a abrir un archivo sucio sin instrucciones y a buscar qué está mal. Descubrir la suciedad solo es la mitad del oficio.}

\milcestacion{milcVerde}{Estación 1 de 3 · Escucha}

\begin{softbox}{milcVerde}{milcGris}{Escena para escuchar}
\textbf{Actividad 1 · IDENTIFICA --- Descubre la suciedad} (15 min · individual). (1) Abre el archivo sucio que reparte tu docente, sin instrucciones sobre qué buscar. (2) Cuenta cuántas filas parecen repetidas. (3) Busca cuántos formatos distintos de fecha aparecen. (4) Cuenta las celdas vacías y mira si se concentran en alguna columna. (5) Busca nombres escritos de dos maneras que en realidad son el mismo. (6) Anota cuál de los cuatro problemas te costó más ver.
\end{softbox}

{\sffamily\bfseries\fontsize{8}{10}\selectfont\textcolor{milcNegro!55}{MARCA CUANDO LO HAYAS HECHO}}
\milccheckrow{Conté duplicados, formatos de fecha y celdas vacías.}
\milccheckrow{Encontré al menos un nombre escrito de dos maneras.}
\milccheckrow{Anoté cuál problema me costó más ver.}

\begin{infoband}{milcVerde}


\textbf{Lo que va al cuaderno (Actividad 1 · IDENTIFICA).} \textbf{Título:} «Actividad 1 --- Descubre la suciedad». \textbf{Formato:} tabla de 4 filas y 3 columnas (tipo de suciedad / cuántos encontré / dónde se concentran), con marca en el que costó más ver. \textbf{Extensión:} un tercio de página. \textbf{Sabes que terminaste cuando} los cuatro tipos tienen un número, aunque alguno sea cero.
\end{infoband}

\puente{Ya la encontraste. Ahora vas a ordenarla en cuatro tipos y a decidir el tratamiento de cada uno antes de tocar nada.}

\milcestacion{milcTurquesa}{Estación 2 de 3 · Entender}

\begin{softbox}{milcTurquesa}{milcGris}{Lo que vas a entender}
\textbf{Actividad 2 · EXPLICA --- Los cuatro tipos y su tratamiento} (20 min · parejas). (1) Con tu pareja, escriban los cuatro tipos con una frase propia cada uno. (2) Escriban el tratamiento de cada uno antes de tocar el archivo. (3) Discutan los vacíos: decidan para cada columna si se completan, se marcan o se descarta la fila, y escriban por qué. (4) Anoten qué se perdería si borraran las filas incompletas sin avisar.
\end{softbox}

\milcpaso{1}{milcTurquesa}{\textbf{Duplicados.} Filas que repiten información. Se elimina la copia y se conserva el original, siempre con nota. Cuidado con los falsos duplicados: dos personas distintas pueden llamarse igual.}
\milcpaso{2}{milcTurquesa}{\textbf{Formatos y errores de escritura.} La misma fecha escrita de tres modos, o un nombre tecleado de dos maneras, parten en dos lo que era uno. Al normalizar suele bajar el número de categorías, y esa bajada es la señal de que funcionó.}
\milcpaso{3}{milcTurquesa}{\textbf{Los vacíos son una decisión, no un accidente.} Se completan, se marcan con un valor explícito o se descarta la fila. Las tres opciones son válidas; lo que no vale es no elegir y dejar que otro adivine.}
\milcpaso{4}{milcTurquesa}{\textbf{Identificar, decidir, documentar.} En ese orden. Nunca se corrige sin haber decidido, y nunca se decide sin dejar constancia. El original se conserva y la copia se limpia.}

\begin{drawbox}{milcTurquesa}{Los cuatro tipos de suciedad}
\textbf{Duplicados:} eliminar la copia, conservar el original. \\[1mm] \textbf{Formatos:} elegir uno y normalizar todo. \\[1mm] \textbf{Vacíos:} completar, marcar o descartar, y decirlo. \\[1mm] \textbf{Escritura:} unificar lo que era lo mismo. \\[1mm] \textbf{Regla:} original intacto, copia limpia, bitácora abierta.
\end{drawbox}

\begin{softbox}{milcMostaza}{milcGris}{Errores comunes y su corrección}
(1) \textbf{Limpiar sobre el original}: sin copia no hay vuelta atrás. (2) \textbf{Borrar filas incompletas en bloque}: se van casos reales sin que nadie sepa cuántos. (3) \textbf{Eliminar duplicados sin mirar}: dos personas pueden llamarse igual. (4) \textbf{Corregir sin documentar}: al día siguiente nadie recuerda por qué ese dato cambió. (5) \textbf{Atender todo a la vez}: se resuelve un tipo de problema por pasada, o se pierde el hilo.
\end{softbox}

\begin{infoband}{milcTurquesa}


\textbf{Lo que va al cuaderno (Actividad 2 · EXPLICA).} \textbf{Título:} «Actividad 2 --- Los cuatro tipos y su tratamiento». \textbf{Formato:} los cuatro tipos con frase propia, su tratamiento decidido antes de tocar nada y la decisión sobre los vacíos columna por columna. \textbf{Extensión:} media página. \textbf{Sabes que terminaste cuando} cada columna con vacíos tiene su decisión escrita.
\end{infoband}

\puente{Tienes el método. Toca limpiar de verdad, sobre una copia y con la bitácora abierta desde el primer cambio.}

\milcestacion{milcMagenta}{Estación 3 de 3 · Hacer}

\begin{softbox}{milcMagenta}{milcGris}{Cómo vas a construir tu producto}
\textbf{Actividad 3 · APLICA --- Limpia con bitácora} (30 min · individual). (1) Duplica el archivo y trabaja solo sobre la copia. (2) Abre la bitácora antes de tocar el primer dato. (3) Resuelve los duplicados y anota cuántos eran y cuáles conservaste. (4) Normaliza los formatos y anota cuál elegiste. (5) Resuelve los vacíos según lo decidido y anota la regla que aplicaste. (6) Unifica los errores de escritura y pásale la bitácora a un compañero para que repita tu limpieza.
\end{softbox}

\milcpaso{1}{milcMagenta}{\textbf{Tu entrega tiene tres piezas.} El original intacto. La copia limpia. Y la bitácora, con al menos cuatro entradas, una por tipo de suciedad.}
\milcpaso{2}{milcMagenta}{\textbf{Cada entrada responde cuatro cosas.} Qué encontraste, qué decidiste, por qué y a qué celdas afectó. Sin la tercera no se puede discutir la decisión; sin la cuarta no se puede deshacer.}
\milcpaso{3}{milcMagenta}{\textbf{Un tipo de problema por pasada.} Primero duplicados, después formatos, después vacíos, después escritura. Mezclarlos hace que un cambio tape a otro y que la bitácora deje de cuadrar.}
\milcpaso{4}{milcMagenta}{\textbf{La prueba es que otro pueda repetirla.} Si tu compañero, con la bitácora en la mano y el archivo original, llega al mismo resultado, está bien documentada. Si tiene que preguntarte algo, falta una línea.}

\begin{drawbox}{milcMagenta}{Antes de entregar}
\checkbox El archivo original quedó intacto. \\[1mm] \checkbox La limpieza se hizo sobre una copia. \\[1mm] \checkbox Bitácora con al menos cuatro entradas, una por tipo. \\[1mm] \checkbox Cada entrada dice qué, qué decidí, por qué y qué celdas. \\[1mm] \checkbox Las columnas y el sentido del original se conservan. \\[1mm] \checkbox Un compañero repitió la limpieza leyendo solo la bitácora.
\end{drawbox}

\begin{softbox}{milcMagenta}{milcGris}{Plantilla de trabajo}
\textbf{Qué encontré.} \textit{Tipo de suciedad y cuántos casos.} \\[1mm] \textbf{Qué decidí.} \textit{El tratamiento elegido.} \\[1mm] \textbf{Por qué.} \textit{La razón, en una línea.} \\[1mm] \textbf{Qué celdas.} \textit{Columnas y filas afectadas.} \\[1mm] \textbf{Prueba.} \textit{Quién repitió la limpieza y si llegó al mismo resultado.}
\end{softbox}

\begin{infoband}{milcMagenta}


\textbf{Lo que va al cuaderno (Actividad 3 · APLICA).} \textbf{Título:} «Actividad 3 --- Limpia con bitácora». \textbf{Formato:} las cuatro entradas de la bitácora completas y la nota de quién repitió la limpieza y con qué resultado. \textbf{Extensión:} media página. \textbf{Sabes que terminaste cuando} un compañero llegó al mismo resultado con la bitácora sola.
\end{infoband}

\puente{Lo que entregas hoy: el archivo limpio y su bitácora. La prueba de calidad: un compañero repite tu limpieza leyendo solo la bitácora.}

\milcestacion{milcMostaza}{Producto de la sesión}
\textbf{Conjunto de datos limpio sobre copia + bitácora reproducible de cuatro entradas}.
\begin{softbox}{milcMostaza}{milcGris}{Criterios de calidad}
(1) El archivo original quedó intacto y la limpieza se hizo sobre una copia. (2) La bitácora tiene al menos cuatro entradas, una por tipo de suciedad. (3) Cada entrada dice qué encontraste, qué decidiste, por qué y a qué celdas afectó. (4) El conjunto limpio conserva las columnas y el sentido del original. (5) Un compañero repitió la limpieza leyendo solo la bitácora.
\end{softbox}

\puente{Antes de cerrar, mira lo que hiciste desde cinco lados. Borrar sin dejar constancia no es limpiar: es hacer desaparecer.}

\milcestacion{milcVino}{Evaluación liberadora · cinco dimensiones}

\begin{softbox}{milcVino}{milcGris}{Sentido de la evaluación}
Evaluar no es solo revisar una nota. Es reconocer qué aprendí, qué cambió en mí, qué emociones manejé, cómo me relaciono con la ciudad y la comunidad, y qué le dejaré a quien viene detrás.
\end{softbox}

\textbf{Mi reflexión liberadora en cinco dimensiones.} Completa cada frase iniciada:

\small
\begin{tabularx}{\textwidth}{>{\bfseries\RaggedRight\arraybackslash}p{3.4cm}Y>{\RaggedRight\arraybackslash}p{4.6cm}}
\rowcolor{milcVino}\color{white}Dimensión & \color{white}\bfseries Mi respuesta (frame starter) & \color{white}\bfseries Algo que descubrí\\
1. Personal & Lo que más cambió en mí fue \linead{6.5cm} & \linead{4.2cm}\\
2. Emocional & La emoción más fuerte fue \linead{2.5cm} y la regulé \linead{3cm} & \linead{4.2cm}\\
3. Ciudadana & En democracia esto importa porque \linead{6cm} & \linead{4.2cm}\\
4. Local (Cartago) & En mi barrio sería útil para \linead{6.5cm} & \linead{4.2cm}\\
5. Intergeneracional & A mi abuela/abuelo le diría \linead{6.5cm} & \linead{4.2cm}\\
\end{tabularx}
\normalsize

\begin{infoband}{milcVino}
\textbf{Para la próxima sustentación.} Vuelve a esta tabla en seis meses. Verás que las respuestas cambian — no porque lo aprendido cambie, sino porque tú cambias. La evaluación liberadora es una conversación contigo a través del tiempo.
\end{infoband}


\puente{Para terminar, tres citas verificadas sobre lo que calla un registro: Dussel sobre la lengua que se comprende por lo que calla, Séneca sobre la demora como remedio, Floridi sobre cribar los datos valiosos.}

\milcestacion{milcGradeDark}{Triángulo de pensamiento · cierre filosófico}

\begin{softbox}{milcVino}{milcGris}{Dussel · lente del nosotros}
\textit{«Hay lengua cotidiana (la de todos los días), lengua de culturas ilustradas, de cultura de masas, lengua de cultura popular… lengua política (que se comprende no por lo que dice sino por lo que calla, contra quién lo dice, cuándo y por qué…).»}

{\footnotesize\itshape\color{milcNegro!70}--- Enrique Dussel · Filosofía de la liberación (1977), §4.2.4.4}

Un conjunto de datos también se comprende por lo que calla. Las celdas vacías no son ruido: dicen a quién no se le preguntó, o quién no quiso contestar. Anotarlas es lo que impide que ese silencio se pierda.

\textbf{¿Qué calla mi tabla, y estoy documentando ese silencio o borrándolo?}
\end{softbox}
\linead{\linewidth}\\[1mm]
\linead{\linewidth}

\begin{softbox}{milcMostaza}{milcGris}{Estoicismo · Séneca · Sobre la ira, II, 29 (c. 45 d.C.) · cuidado interior}
\textit{«El mayor remedio para la ira es la demora. (trad. propia)»}

{\footnotesize\itshape\color{milcNegro!70}--- Séneca · Sobre la ira, II, 29 (c. 45 d.C.)}

Con la limpieza pasa lo mismo. La tentación es borrar de una todas las filas incompletas y seguir. Esperar el minuto que toma decidir qué hacer con ellas es lo que separa limpiar de hacer desaparecer.

\textbf{¿Qué borré de un tirón alguna vez y después necesité de vuelta?}
\end{softbox}
\linead{\linewidth}\\[1mm]
\linead{\linewidth}

\begin{softbox}{milcTurquesa}{milcGris}{Floridi · lente de la infoesfera}
\textit{«Necesitamos más y mejores técnicas y tecnologías para ver los pequeños patrones en los datos, pero necesitamos más y mejor epistemología para cribar los valiosos. (trad. propia)»}

{\footnotesize\itshape\color{milcNegro!70}--- Luciano Floridi · Big data and their epistemological challenge (2012)}

Las herramientas encuentran duplicados en un segundo. Decidir si dos filas iguales son un error o dos personas que se llaman igual no lo resuelve ninguna herramienta: lo resuelve quien conoce los datos.

\textbf{¿Qué decisión de mi limpieza no habría podido tomar un programa por mí?}
\end{softbox}
\linead{\linewidth}\\[1mm]
\linead{\linewidth}

\begin{infoband}{milcNegro}
\textbf{Nota para el docente.} Las tres son citas textuales y llevan su referencia debajo. Puedes remitir a la obra en clase.
\end{infoband}

\milcestacion{milcVino}{Mi autoevaluación · marca una casilla por fila}

{\small
\setlength{\extrarowheight}{2.2pt}
\begin{tabularx}{\textwidth}{>{\RaggedRight\arraybackslash\bfseries}p{3.0cm}YYY}
\rowcolor{milcVino}\color{white}\bfseries Criterio & \color{white}\bfseries Lo logré & \color{white}\bfseries En proceso & \color{white}\bfseries Necesito apoyo\\[.6mm]
Comprensión del tema & \milcopcion{Explico las ideas con mis palabras.} & \milcopcion{Reconozco las ideas, falta precisión.} & \milcopcion{Necesito repasar lo básico.}\\[.8mm]
\rowcolor{milcGris}Pensamiento computacional & \milcopcion{Usé los pilares al armar mi producto.} & \milcopcion{Usé algunos pilares.} & \milcopcion{Necesito releer la estación 2.}\\[.8mm]
Aplicación y producto & \milcopcion{Mi producto cumple los criterios.} & \milcopcion{Está armado, con ajustes pendientes.} & \milcopcion{Aún está en construcción.}\\[.8mm]
\rowcolor{milcGris}Reflexión 5 dimensiones & \milcopcion{Respondí cada una con honestidad.} & \milcopcion{Respondí algunas con detalle.} & \milcopcion{Mis respuestas son muy generales.}\\[.8mm]
Triángulo de pensamiento & \milcopcion{Conecté las 3 voces con mi trabajo.} & \milcopcion{Conecté al menos una.} & \milcopcion{No las relaciono con mi proceso.}\\[.8mm]
\end{tabularx}}

\vspace{3mm}
\textbf{Hoy entendí que:}\\[1mm]
\linead{\linewidth}\\[1mm]
\linead{\linewidth}

\textbf{Mi compromiso para la próxima sesión es:}\\[1mm]
\linead{\linewidth}\\[1mm]
\linead{\linewidth}

\begin{softbox}{milcMostaza}{milcGris}{Cita de despedida · Marco Aurelio}
\textit{``El obstáculo en el camino se vuelve el camino. Lo que estorba al camino se vuelve camino.''}\\[1mm]
Cada error de hoy, bien observado, se convierte en pieza del camino que pisarás mañana.
\end{softbox}

\small
\begin{tabularx}{\textwidth}{>{\bfseries}p{3.6cm}Y}
\rowcolor{milcGradeDark!12}Nombre completo: & \linead{10cm}\\
Fecha: & \linead{4cm} \quad Curso/grupo: \linead{4cm}\\
Mi firma: & \linead{9cm}\\
\end{tabularx}
\normalsize

\begin{infoband}{milcGradeDark}
\textbf{Para la próxima sesión.} Llega con la copia limpia, el original intacto y la bitácora. Pídele a un compañero que repita tu limpieza y anota dónde se atascó. En la sesión 8 vas a auditar gráficos que engañan, y un dato mal limpiado es la mitad de un gráfico falso.
\end{infoband}

\end{document}
